\documentclass{article}%
\usepackage[T1]{fontenc}%
\usepackage[utf8]{inputenc}%
\usepackage{lmodern}%
\usepackage{textcomp}%
\usepackage{lastpage}%
\usepackage{geometry}%
\geometry{margin=1in,headheight=14pt,headsep=25pt}%
\usepackage{hyperref}%
\usepackage{enumitem}%
\usepackage{fancyhdr}%
\usepackage{xcolor}%
%

            \hypersetup{
                colorlinks=true,
                linkcolor=blue,
                filecolor=magenta,
                urlcolor=blue,
            }
            
            \pagestyle{fancy}
            \fancyhf{}
            \rhead{Generated on \today}
            \lhead{Course Summary}
            \cfoot{\thepage}
        %
\title{Course Summary}%
\author{Generated by ScribeLec}%
\date{\today}%
%
\begin{document}%
\normalsize%
\maketitle%
\tableofcontents\newpage%
\section{Key Terms}%
\label{sec:KeyTerms}%
\begin{itemize}[[leftmargin=*]]%
\item \href{https://www.math.purdue.edu/~yipn/421/2024-08-27-ExSimplex.pdf#page=1}{\textbf{Simplex Method}}: An iterative algorithm used to solve linear programming problems by moving from one vertex of the feasible region to another until an optimal solution is found.<L[2024-08-27-ExSimplex] 1><L[2024-08-27-ExSimplex] 6>%
\item \href{https://www.math.purdue.edu/~yipn/421/2024-08-27-ExSimplex.pdf#page=1}{\textbf{Slack Variable}}: A non-negative variable added to a less-than-or-equal-to inequality constraint to transform it into an equality constraint.<L[2024-08-27-ExSimplex] 1>%
\item \href{https://www.math.purdue.edu/~yipn/421/2024-08-27-ExSimplex.pdf#page=2}{\textbf{Feasible Region}}: The set of all points that satisfy all the constraints of a linear programming problem.<L[2024-08-27-ExSimplex] 2><L[2024-08-27-ExSimplex] 4><L[2024-08-27-ExSimplex] 5><L[2024-08-27-ExSimplex] 16><L[2024-08-27-ExSimplex] 30>%
\item \href{https://www.math.purdue.edu/~yipn/421/2024-08-27-ExSimplex.pdf#page=9}{\textbf{Basic Variable}}: In the Simplex method, a variable that is expressed in terms of non-basic variables in a dictionary.<L[2024-08-27-ExSimplex] 9><L[2024-08-27-ExSimplex] 10><L[2024-08-27-ExSimplex] 12><L[2024-08-27-ExSimplex] 13><L[2024-08-27-ExSimplex] 14><L[2024-08-27-ExSimplex] 15><L[2024-08-27-ExSimplex] 26><L[2024-08-27-ExSimplex] 27><L[2024-08-27-ExSimplex] 28>%
\item \href{https://www.math.purdue.edu/~yipn/421/2024-08-27-ExSimplex.pdf#page=6}{\textbf{Non{-}basic Variable}}: In the Simplex method, a variable set to zero in a given iteration of the algorithm.<L[2024-08-27-ExSimplex] 6><L[2024-08-27-ExSimplex] 7><L[2024-08-27-ExSimplex] 26><L[2024-08-27-ExSimplex] 27><L[2024-08-27-ExSimplex] 28>%
\item \href{https://www.math.purdue.edu/~yipn/421/2024-08-27-ExSimplex.pdf#page=9}{\textbf{Pivot Operation}}: The process of exchanging a basic and a non-basic variable in the Simplex method to move to a new vertex of the feasible region.<L[2024-08-27-ExSimplex] 9><L[2024-08-27-ExSimplex] 10><L[2024-08-27-ExSimplex] 12><L[2024-08-27-ExSimplex] 13><L[2024-08-27-ExSimplex] 14><L[2024-08-27-ExSimplex] 15><L[2024-08-27-ExSimplex] 27><L[2024-08-27-ExSimplex] 28>%
\item \href{https://www.math.purdue.edu/~yipn/421/2024-08-27-ExSimplex.pdf#page=26}{\textbf{Dictionary}}: A representation of a linear programming problem where basic variables are expressed as functions of non-basic variables.<L[2024-08-27-ExSimplex] 26><L[2024-08-27-ExSimplex] 27><L[2024-08-27-ExSimplex] 28><L[2024-08-27-ExSimplex] 29>%
\item \href{https://www.math.purdue.edu/~yipn/421/2024-08-27-ExSimplex.pdf#page=32}{\textbf{Auxiliary Problem}}: A modified linear programming problem introduced to find an initial feasible solution when the origin is not feasible in the original problem.<L[2024-08-27-ExSimplex] 32><L[2024-08-27-ExSimplex] 33><L[2024-08-27-ExSimplex] 34><L[2024-08-27-ExSimplex] 35><L[2024-08-27-ExSimplex] 36><L[2024-08-27-ExSimplex] 37><L[2024-08-27-ExSimplex] 38><L[2024-08-27-ExSimplex] 39><L[2024-08-27-ExSimplex] 40><L[2024-08-27-ExSimplex] 41><L[2024-08-27-ExSimplex] 42>%
\end{itemize}%
\newpage

%
\section{Problem Types}%
\label{sec:ProblemTypes}%
\begin{itemize}[[leftmargin=*]]%
\item \href{https://www.math.purdue.edu/~yipn/421/2024-08-27-ExSimplex.pdf#page=32}{\textbf{Finding an Initial Feasible Solution}}: The problem of finding a feasible solution to start the Simplex method, especially when the origin is not feasible. An auxiliary problem is introduced to find a feasible starting point. <L[2024-08-27-ExSimplex] 32><L[2024-08-27-ExSimplex] 33><L[2024-08-27-ExSimplex] 34><L[2024-08-27-ExSimplex] 35><L[2024-08-27-ExSimplex] 36><L[2024-08-27-ExSimplex] 37><L[2024-08-27-ExSimplex] 38><L[2024-08-27-ExSimplex] 39><L[2024-08-27-ExSimplex] 40><L[2024-08-27-ExSimplex] 41><L[2024-08-27-ExSimplex] 42>%
\item \href{https://www.math.purdue.edu/~yipn/421/2024-08-27-ExSimplex.pdf#page=30}{\textbf{Handling Infeasible Origins}}: Addressing situations where the origin (0,0) or (0,0,0) does not satisfy the constraints of the linear programming problem, requiring a different starting point for the Simplex method. <L[2024-08-27-ExSimplex] 30><L[2024-08-27-ExSimplex] 32><L[2024-08-27-ExSimplex] 33><L[2024-08-27-ExSimplex] 34><L[2024-08-27-ExSimplex] 35><L[2024-08-27-ExSimplex] 36><L[2024-08-27-ExSimplex] 37><L[2024-08-27-ExSimplex] 38><L[2024-08-27-ExSimplex] 39><L[2024-08-27-ExSimplex] 40><L[2024-08-27-ExSimplex] 41><L[2024-08-27-ExSimplex] 42><L[2024-08-27-ExSimplex] 43>%
\item \href{https://www.math.purdue.edu/~yipn/421/2024-08-27-ExSimplex.pdf#page=2}{\textbf{Graphical Representation of Linear Inequalities and Feasible Regions}}: Visualizing the constraints of a linear programming problem in two dimensions to identify the feasible region (the set of points satisfying all constraints). This aids in understanding the Simplex method's iterative process of moving between vertices. <L[2024-08-27-ExSimplex] 2><L[2024-08-27-ExSimplex] 3><L[2024-08-27-ExSimplex] 4><L[2024-08-27-ExSimplex] 5><L[2024-08-27-ExSimplex] 16>%
\item \href{https://www.math.purdue.edu/~yipn/421/2024-08-27-ExSimplex.pdf#page=6}{\textbf{Simplex Method Iteration and Pivot Operations}}: The iterative process of the Simplex method, involving moving from one vertex of the feasible region to another by exchanging basic and non-basic variables (pivot operation) to improve the objective function value. <L[2024-08-27-ExSimplex] 6><L[2024-08-27-ExSimplex] 7><L[2024-08-27-ExSimplex] 8><L[2024-08-27-ExSimplex] 9><L[2024-08-27-ExSimplex] 10><L[2024-08-27-ExSimplex] 11><L[2024-08-27-ExSimplex] 12><L[2024-08-27-ExSimplex] 13><L[2024-08-27-ExSimplex] 14><L[2024-08-27-ExSimplex] 15><L[2024-08-27-ExSimplex] 17><L[2024-08-27-ExSimplex] 18><L[2024-08-27-ExSimplex] 19><L[2024-08-27-ExSimplex] 20><L[2024-08-27-ExSimplex] 21><L[2024-08-27-ExSimplex] 22><L[2024-08-27-ExSimplex] 23><L[2024-08-27-ExSimplex] 24><L[2024-08-27-ExSimplex] 25>%
\item \href{https://www.math.purdue.edu/~yipn/421/2024-08-27-ExSimplex.pdf#page=26}{\textbf{Using Dictionaries to Represent Linear Programs}}: Organizing the variables of a linear program into basic and non-basic sets, represented by a dictionary of equations, to facilitate the Simplex method's iterative process. <L[2024-08-27-ExSimplex] 26><L[2024-08-27-ExSimplex] 27><L[2024-08-27-ExSimplex] 28><L[2024-08-27-ExSimplex] 29>%
\item \href{https://www.math.purdue.edu/~yipn/421/2024-08-27-ExSimplex.pdf#page=17}{\textbf{Solving Linear Programs with Three Variables}}: Extending the Simplex method to linear programming problems involving three variables, requiring a more complex iterative process. <L[2024-08-27-ExSimplex] 17><L[2024-08-27-ExSimplex] 18><L[2024-08-27-ExSimplex] 19><L[2024-08-27-ExSimplex] 20><L[2024-08-27-ExSimplex] 21><L[2024-08-27-ExSimplex] 22><L[2024-08-27-ExSimplex] 23><L[2024-08-27-ExSimplex] 24><L[2024-08-27-ExSimplex] 25>%
\end{itemize}%
\newpage

%
\section{Algorithm Solutions}%
\label{sec:AlgorithmSolutions}%
\begin{itemize}[[leftmargin=*]]%
\item \href{https://www.math.purdue.edu/~yipn/421/2024-08-27-ExSimplex.pdf#page=6}{\textbf{Simplex Method}}: Iteratively move from one vertex of the feasible region to another by setting non-basic variables to zero and choosing a new set of (non)basic variables to improve the objective function, until the optimal solution is found.<L[2024-08-27-ExSimplex] 6><L[2024-08-27-ExSimplex] 7><L[2024-08-27-ExSimplex] 8><L[2024-08-27-ExSimplex] 9><L[2024-08-27-ExSimplex] 10><L[2024-08-27-ExSimplex] 11><L[2024-08-27-ExSimplex] 12><L[2024-08-27-ExSimplex] 13><L[2024-08-27-ExSimplex] 14><L[2024-08-27-ExSimplex] 15><L[2024-08-27-ExSimplex] 16><L[2024-08-27-ExSimplex] 17><L[2024-08-27-ExSimplex] 18><L[2024-08-27-ExSimplex] 19><L[2024-08-27-ExSimplex] 20><L[2024-08-27-ExSimplex] 21><L[2024-08-27-ExSimplex] 22><L[2024-08-27-ExSimplex] 23><L[2024-08-27-ExSimplex] 24><L[2024-08-27-ExSimplex] 25>%
\item \href{https://www.math.purdue.edu/~yipn/421/2024-08-27-ExSimplex.pdf#page=32}{\textbf{Auxiliary Problem}}: Introduce a new variable  $x_0$ and modify constraints to create a new problem whose feasible region includes the origin. Solve this problem to find a feasible solution for the original problem, then proceed with the Simplex method.<L[2024-08-27-ExSimplex] 32><L[2024-08-27-ExSimplex] 33><L[2024-08-27-ExSimplex] 34><L[2024-08-27-ExSimplex] 35><L[2024-08-27-ExSimplex] 36><L[2024-08-27-ExSimplex] 37><L[2024-08-27-ExSimplex] 38><L[2024-08-27-ExSimplex] 39><L[2024-08-27-ExSimplex] 40><L[2024-08-27-ExSimplex] 41><L[2024-08-27-ExSimplex] 42><L[2024-08-27-ExSimplex] 43>%
\end{itemize}%
\newpage

%
\end{document}